% !TEX root = ../main.tex
\chapter{\cmm{}: entrada, saída e biblioteca padrão}
\label{cap:cmm-io-stdlib}

\section{Entrada e saída}
\label{sec:io}

O processador conversa com o mundo por portas numeradas, criadas pelas diretivas \code{\#NUIOIN} e \code{\#NUIOOU}, e todo o tráfego passa por quatro funções intrínsecas: \code{in(p)} lê um inteiro da porta de entrada \code{p}; \code{fin(p)} lê convertendo para \code{float}; \code{out(p, expr);} escreve o valor da expressão na porta de saída \code{p}; e \code{fout(p, expr);} escreve em formato \code{float}.

O número da porta é um literal inteiro, validado contra as diretivas: usar \code{in(2)} num processador com \code{\#NUIOIN 1} é erro de compilação. No hardware, \code{in()} é uma leitura com \textit{handshake}, na qual o processador sinaliza \code{req\_in} e aguarda o dado, e \code{out()} pulsa \code{out\_en} com o dado estável no barramento, o suficiente para acoplar FIFOs e conversores. No nosso \code{media_movel}, \code{x[0] = in(0);} lê a amostra e \code{out(0, soma >> 2);} publica a média: uma porta de cada lado basta.

\section{Funções especiais}

Cinco funções existem para economizar hardware, cada uma evitando um bloco caro da ULA ou uma sequência de instruções. A \code{norm(x)} divide pelo valor da diretiva \code{\#NUGAIN}, uma potência de dois, e entrega a normalização sem instanciar o divisor completo; funciona só com inteiros. A \code{pset(x)} devolve \code{x} quando positivo e zero quando negativo, o equivalente de \code{if (x<0) x=0;} numa única instrução. A \code{abs(x)} é o valor absoluto e, sobre um complexo, a magnitude. A \code{sign(x, y)} devolve \code{y} com o sinal de \code{x}. E a \code{copy(x, y);} copia os bits de \code{x} em \code{y} sem conversão nem checagem de tipo, uma reinterpretação crua, a ser usada com consciência.

\section{Funções não lineares e de arredondamento}

A biblioteca cobre raiz e trigonometria (\code{sqrt}, \code{sin}, \code{cos}, \code{tan}, \code{atan}), as hiperbólicas (\code{sinh}, \code{cosh}, \code{tanh}), exponencial e logaritmo (\code{exp}, \code{log}, \code{pow}) e os arredondamentos \code{floor}, \code{ceil} e \code{round}, que retornam \code{float}, com o \code{round} arredondando o meio para longe do zero.

Essas funções não são circuitos dedicados: são macros de \textit{assembly} otimizadas, como o método de Newton na \code{sqrt} e séries nas trigonométricas, injetadas no programa quando usadas. Elas custam instruções e ciclos, não blocos de ULA, exceto pelas operações de ponto flutuante que elas próprias empregam. A \code{pow(x, y)} merece nota: com expoente inteiro constante, o compilador gera a sequência mínima de multiplicações; com expoente inteiro variável, um laço em tempo de execução; nos demais casos, a identidade $x^y = e^{y \ln x}$.

\section{Funções para complexos}

Seis funções operam sobre o tipo \code{comp}: \code{real(z)} e \code{imag(z)} extraem as partes; \code{fase(z)} devolve o argumento; \code{mod2(z)} devolve a magnitude ao quadrado, $a^2+b^2$, mais barata que \code{abs} por dispensar a raiz; \code{complex(re, im)} monta um complexo a partir de dois \code{float}; e \code{conj(z)} devolve o conjugado.

\begin{lstlisting}[style=cmm]
comp x; comp y; comp r;
x = 1.0 + 2.0i;
y = 3.0 + 4.0i;
r = x * y;
fout(0, real(r));
fout(0, imag(r));
\end{lstlisting}

\begin{aviso}
Nem tudo se aplica a \code{comp}: \code{sqrt}, \code{exp}, \code{log}, \code{pow}, \code{sign} e \code{pset} sobre complexos são erro de compilação, assim como o incremento \code{z++}. O módulo \code{\%} e a \code{norm()} são exclusivos de inteiros.
\end{aviso}

\section{Álgebra linear em notação de Dirac}
\label{sec:dirac}

O recurso mais idiossincrático da \cmm{} é a escrita de operações vetoriais e matriciais em notação \textit{bra-ket}, familiar a quem vem da física. Sendo \code{a} e \code{b} vetores (\textit{arrays}) e \code{M} uma matriz, o produto interno é a expressão \texttt{<a|b>}; o comando \texttt{a \# |0>;} zera o vetor; \texttt{a \# |M|b>;} calcula o produto matriz-vetor; \texttt{a \# c|b>;} escala um vetor por uma constante; \texttt{A \# |a><b|;} forma o produto externo; e \texttt{out(p, c|a>);} emite o vetor escalado por uma porta de saída.

Essas construções expandem para os laços e instruções indexadas adequados, e são a forma idiomática de escrever filtros FIR, correlações e projeções em \cmm{}. Um FIR direto, por exemplo, é uma linha, \texttt{y = <h|x>;}, com \code{h} guardando os coeficientes, inicializados por arquivo, e \code{x} a janela de amostras.
